% BHU PYQ -- Transcribed & Corrected
% Paper: CHEMV-31 / CHEMV31: Industrial Chemistry - I
% Exam: B.Sc. (Hons.) Semester III Examination, 2025-26 (NEP)
% Category: Minor / Multidisciplinary Course (MD)
% Department: Chemistry

\documentclass[11pt,a4paper]{article}
\usepackage[a4paper,top=2.5cm,bottom=2.5cm,left=2.8cm,right=2.8cm,headheight=14pt]{geometry}
\usepackage{amsmath,amssymb}
\usepackage[T1]{fontenc}
\usepackage[utf8]{inputenc}
\usepackage{lmodern,microtype,fancyhdr,enumitem,hyperref,lastpage,parskip}

\hypersetup{
    colorlinks=true,
    urlcolor=black,
    linkcolor=black,
    pdftitle={CHEMV31: Industrial Chemistry I -- BHU Sem III 2025-26 (NEP)}
}

\pagestyle{fancy}
\fancyhf{}
\renewcommand{\headrulewidth}{0.4pt}
\renewcommand{\footrulewidth}{0.4pt}
\fancyhead[L]{\small\textit{Banaras Hindu University}}
\fancyhead[R]{\small\textit{CHEMV-31: Industrial Chemistry - I}}
\fancyfoot[C]{\small Page \thepage\ of \pageref{LastPage}}

\newlist{parts}{enumerate}{2}
\setlist[parts,1]{label=(\alph*),leftmargin=*,itemsep=4pt,topsep=2pt}
\setlist[parts,2]{label=(\roman*),leftmargin=*,itemsep=2pt,topsep=2pt}
\newcommand{\pts}[1]{\hfill\textit{[#1]}}

\begin{document}

\begin{center}
  {\large\bfseries BANARAS HINDU UNIVERSITY}\\[2pt]
  {\normalsize B.Sc. (Hons.) --- Semester III Examination, 2025-26 (NEP)}\\[6pt]
  \rule{0.8\linewidth}{0.5pt}\\[6pt]
  {\large\bfseries CHEMISTRY}\\[4pt]
  {\large\bfseries Paper Code: CHEMV-31 : Industrial Chemistry --- I}\\[4pt]
  {\normalsize Time: 2:30 Hours\hfill Maximum Marks: 70}\\[2pt]
  {\small REF NO. CECONF/N/2025-26/058}
\end{center}

\rule{\linewidth}{0.4pt}

\smallskip
\noindent\textit{Instructions: Answer 	extbf{five} questions in all including Question No.~1 which is compulsory. All questions carry equal marks (14 marks each).

\bigskip

\noindent\textbf{Question 1.} Short Answer Questions (Compulsory):\pts{$2 \times 7 = 14$}
\begin{parts}
  \item Discuss the INCI naming system used in cosmetic labeling.\pts{2}
  \item Explain the role of builders in detergent formulations.\pts{2}
  \item What is the significance of the INCI name on cosmetic labels?\pts{2}
  \item Who launched the first vegetarian soap in India?\pts{2}
  \item Differentiate between soft soaps and hard soaps.\pts{2}
  \item What is a superfatted soap?\pts{2}
  \item Why is brine (sodium chloride solution) used in the soap-making process (\textit{salting out})?\pts{2}
\end{parts}

\medskip\rule{\linewidth}{0.2pt}\medskip

\noindent\textbf{Question 2.}
\begin{parts}
  \item Classify the fatty acids used in soap preparation and explain their role in controlling soap physical and chemical properties with suitable examples.\pts{7}
  \item What are the main ingredients of a typical skin-cleansing product? Mention the specific function of each ingredient.\pts{7}
\end{parts}

\medskip\rule{\linewidth}{0.2pt}\medskip

\noindent\textbf{Question 3.}
\begin{parts}
  \item Write short notes on the following:\pts{7}
  \begin{parts}
    \item Humectants
    \item Transparent soaps
    \item Emollients
    \item Ancient cleaning agents
  \end{parts}
  \item Discuss the concept of Sun Protection Factor (SPF) and mention the main active and inactive ingredients of a sunscreen formulation.\pts{7}
\end{parts}

\medskip\rule{\linewidth}{0.2pt}\medskip

\noindent\textbf{Question 4.}
\begin{parts}
  \item Discuss the environmental impact of non-biodegradable synthetic detergents and explain how biodegradable detergents help to reduce water pollution.\pts{5}
  \item Write short notes on the following:\pts{6}
  \begin{parts}
    \item Surfactants in cosmetic chemistry
    \item Cold process method of soap manufacturing
  \end{parts}
  \item Discuss the mechanism of micelle formation and solubilization in cleaning action.\pts{3}
\end{parts}

\medskip\rule{\linewidth}{0.2pt}\medskip

\noindent\textbf{Question 5.}
\begin{parts}
  \item Differentiate between vanishing creams and cold creams in terms of formulation components and application.\pts{4}
  \item Draw a schematic diagram and briefly describe the working principle of skin hydration measurement using a Corneometer.\pts{4}
  \item Discuss the different levels of packaging (primary, secondary, tertiary) used in cosmetic formulations.\pts{3}
  \item What are anionic, cationic, and non-ionic detergents? Give one chemical example and one specific use of each type.\pts{3}
\end{parts}

\medskip\rule{\linewidth}{0.2pt}\medskip

\noindent\textbf{Question 6.}
\begin{parts}
  \item Write the chemical synthesis of Alcohol Ethoxylate Sulfates (AES) and explain their major industrial applications as detergents.\pts{7}
  \item Mention the various types of lipsticks and the key ingredients used in lipstick formulation.\pts{7}
\end{parts}

\medskip\rule{\linewidth}{0.2pt}\medskip

\noindent\textbf{Question 7.}
\begin{parts}
  \item What are the main layers of epidermis? Describe briefly the structure of the outermost layer of human skin (\textit{stratum corneum}) and its role in skin barrier function and hydration.\pts{4}
  \item Differentiate between soaps and synthetic detergents in terms of chemical composition and effectiveness in hard water.\pts{4}
  \item Briefly discuss the 'hot process' method of soap manufacturing.\pts{3}
  \item Discuss scum formation (\textit{calcium and magnesium soap precipitation}) and how it affects the cleaning efficiency of soaps.\pts{3}
\end{parts}

\vfill
\rule{\linewidth}{0.4pt}
\begin{center}\small
\href{https://bhu-exam.vercel.app/}{Click here to visit our website}\\[4pt]
\href{https://me-aryan.vercel.app/}{Click here to visit our contributor}
\end{center}

\end{document}
